% Options for packages loaded elsewhere
\PassOptionsToPackage{unicode}{hyperref}
\PassOptionsToPackage{hyphens}{url}
\PassOptionsToPackage{dvipsnames,svgnames,x11names}{xcolor}
%
\documentclass[
  letterpaper,
  DIV=11,
  numbers=noendperiod]{scrartcl}

\usepackage{amsmath,amssymb}
\usepackage{iftex}
\ifPDFTeX
  \usepackage[T1]{fontenc}
  \usepackage[utf8]{inputenc}
  \usepackage{textcomp} % provide euro and other symbols
\else % if luatex or xetex
  \usepackage{unicode-math}
  \defaultfontfeatures{Scale=MatchLowercase}
  \defaultfontfeatures[\rmfamily]{Ligatures=TeX,Scale=1}
\fi
\usepackage{lmodern}
\ifPDFTeX\else  
    % xetex/luatex font selection
\fi
% Use upquote if available, for straight quotes in verbatim environments
\IfFileExists{upquote.sty}{\usepackage{upquote}}{}
\IfFileExists{microtype.sty}{% use microtype if available
  \usepackage[]{microtype}
  \UseMicrotypeSet[protrusion]{basicmath} % disable protrusion for tt fonts
}{}
\makeatletter
\@ifundefined{KOMAClassName}{% if non-KOMA class
  \IfFileExists{parskip.sty}{%
    \usepackage{parskip}
  }{% else
    \setlength{\parindent}{0pt}
    \setlength{\parskip}{6pt plus 2pt minus 1pt}}
}{% if KOMA class
  \KOMAoptions{parskip=half}}
\makeatother
\usepackage{xcolor}
\setlength{\emergencystretch}{3em} % prevent overfull lines
\setcounter{secnumdepth}{-\maxdimen} % remove section numbering
% Make \paragraph and \subparagraph free-standing
\makeatletter
\ifx\paragraph\undefined\else
  \let\oldparagraph\paragraph
  \renewcommand{\paragraph}{
    \@ifstar
      \xxxParagraphStar
      \xxxParagraphNoStar
  }
  \newcommand{\xxxParagraphStar}[1]{\oldparagraph*{#1}\mbox{}}
  \newcommand{\xxxParagraphNoStar}[1]{\oldparagraph{#1}\mbox{}}
\fi
\ifx\subparagraph\undefined\else
  \let\oldsubparagraph\subparagraph
  \renewcommand{\subparagraph}{
    \@ifstar
      \xxxSubParagraphStar
      \xxxSubParagraphNoStar
  }
  \newcommand{\xxxSubParagraphStar}[1]{\oldsubparagraph*{#1}\mbox{}}
  \newcommand{\xxxSubParagraphNoStar}[1]{\oldsubparagraph{#1}\mbox{}}
\fi
\makeatother


\providecommand{\tightlist}{%
  \setlength{\itemsep}{0pt}\setlength{\parskip}{0pt}}\usepackage{longtable,booktabs,array}
\usepackage{calc} % for calculating minipage widths
% Correct order of tables after \paragraph or \subparagraph
\usepackage{etoolbox}
\makeatletter
\patchcmd\longtable{\par}{\if@noskipsec\mbox{}\fi\par}{}{}
\makeatother
% Allow footnotes in longtable head/foot
\IfFileExists{footnotehyper.sty}{\usepackage{footnotehyper}}{\usepackage{footnote}}
\makesavenoteenv{longtable}
\usepackage{graphicx}
\makeatletter
\def\maxwidth{\ifdim\Gin@nat@width>\linewidth\linewidth\else\Gin@nat@width\fi}
\def\maxheight{\ifdim\Gin@nat@height>\textheight\textheight\else\Gin@nat@height\fi}
\makeatother
% Scale images if necessary, so that they will not overflow the page
% margins by default, and it is still possible to overwrite the defaults
% using explicit options in \includegraphics[width, height, ...]{}
\setkeys{Gin}{width=\maxwidth,height=\maxheight,keepaspectratio}
% Set default figure placement to htbp
\makeatletter
\def\fps@figure{htbp}
\makeatother

\KOMAoption{captions}{tableheading}
\makeatletter
\@ifpackageloaded{caption}{}{\usepackage{caption}}
\AtBeginDocument{%
\ifdefined\contentsname
  \renewcommand*\contentsname{Table of contents}
\else
  \newcommand\contentsname{Table of contents}
\fi
\ifdefined\listfigurename
  \renewcommand*\listfigurename{List of Figures}
\else
  \newcommand\listfigurename{List of Figures}
\fi
\ifdefined\listtablename
  \renewcommand*\listtablename{List of Tables}
\else
  \newcommand\listtablename{List of Tables}
\fi
\ifdefined\figurename
  \renewcommand*\figurename{Figure}
\else
  \newcommand\figurename{Figure}
\fi
\ifdefined\tablename
  \renewcommand*\tablename{Table}
\else
  \newcommand\tablename{Table}
\fi
}
\@ifpackageloaded{float}{}{\usepackage{float}}
\floatstyle{ruled}
\@ifundefined{c@chapter}{\newfloat{codelisting}{h}{lop}}{\newfloat{codelisting}{h}{lop}[chapter]}
\floatname{codelisting}{Listing}
\newcommand*\listoflistings{\listof{codelisting}{List of Listings}}
\makeatother
\makeatletter
\makeatother
\makeatletter
\@ifpackageloaded{caption}{}{\usepackage{caption}}
\@ifpackageloaded{subcaption}{}{\usepackage{subcaption}}
\makeatother

\ifLuaTeX
  \usepackage{selnolig}  % disable illegal ligatures
\fi
\usepackage{bookmark}

\IfFileExists{xurl.sty}{\usepackage{xurl}}{} % add URL line breaks if available
\urlstyle{same} % disable monospaced font for URLs
\hypersetup{
  pdftitle={Flooding Analysis},
  pdfauthor={Group Members: - Kyle Allie},
  colorlinks=true,
  linkcolor={blue},
  filecolor={Maroon},
  citecolor={Blue},
  urlcolor={Blue},
  pdfcreator={LaTeX via pandoc}}


\title{Flooding Analysis}
\author{}
\date{Invalid Date}

\begin{document}
\maketitle


\section{Front Matter}\label{front-matter}

\begin{itemize}
\tightlist
\item
  \textbf{Course:} Data Analysis STAT 4620/5620 Winter 24-25
\item
  \textbf{GitHub Repository:}
  https://github.com/KyleA2001/STAT4620\_5620-Flooding\_Project.git
\end{itemize}

\section{Abstract}\label{abstract}

Summary in 150 words (Best done last, summarizes the whole project)

\section{Keywords}\label{keywords}

listing out 10 key words max

\section{Introduction}\label{introduction}

Over the past few years, Nova Scotia has faced various floods as a
result of thunderstorms, heavy rain, and extreme weather conditions.
These events severely impacted communities, leading to evacuations,
property damage, infrastructure loss, and loss of life. Certain areas in
Nova Scotia, particularly its coastal regions, are especially vulnerable
to these events. With the increasing frequency and severity of extreme
weather due to climate change, the risk of flooding in Nova Scotia is
expected to grow.

Modeling these rare extreme events is challenging due to their rarity,
yet their societal and economic impacts can be profound and severe. This
research aims to understand the relationships between various climatic
factors and water surge levels at selected locations across Nova Scotia.
The analysis will be followed by predictions of the return periods and
surge levels in the studied regions. Sophisticated and relevant
statistical models will be employed to investigate this research.

This research will address two key questions:

\begin{enumerate}
\def\labelenumi{\arabic{enumi}.}
\item
  What is the relationship between various climatic factors and water
  surge levels across Nova Scotia?
\item
  How can the most significant correlated factors be used to predict
  future flooding catastrophes in Nova Scotia?
\end{enumerate}

Through these investigations, this study aims to contribute valuable
predictions and insights that will aid in understanding and reduce
future flooding risks in the province.

\section{Data Description}\label{data-description}

\begin{table}[ht]
\centering
\begin{tabular}{|l|l|l|l|}
\hline
\textbf{Variables} & \textbf{Units} & \textbf{Description} & \textbf{Frequency} \\ \hline
Date & days & Date when data was recorded & Daily \\ \hline
Location &  & Halifax and Yarmouth &  \\ \hline
Water Level & m (meters) & The height of water above an assumed datum or reference system (any arbitrary local coordinate system) & Hourly \\ \hline
Water Surge & m (meters) & The sudden rise in water level created by both the water level and the tide height (calculated using TideHarmonics R package) & Hourly \\ \hline
Air Temperature & °C (degrees Celsius) & The air temperature & Daily minimum and maximum, Hourly average \\ \hline
Wind speed & km/h (kilometres per hour) & Wind speed & Daily minimum and maximum, Hourly average \\ \hline
Rainfall & mm (millimeters) & The amount of precipitation, not including snow & Real daily value \\ \hline
\end{tabular}
\caption{Data description summary (question 1)}
\end{table}

-- Describe the data being careful to define each variable, including
measurement units if applicable. -- Provide any exploratory
visualisations relevant to the research question

\section{Methods}\label{methods}

-- Provide a brief description of the analytical tools you will used --
Discuss why the method is appropriate in the context of the research
question and the data, including any assumptions required by the method.

\section{Analysis}\label{analysis}

-- Show the code used to actually apply the methods to the data. --
Include model validation.

\section{Results}\label{results}

-- Present and discuss the outputs of your methods. -- Include any
visualisations of the analysis.

\section{Conclusions}\label{conclusions}

-- Discuss the results of the previous section in the context of the
research ques- tion. -- Give an answer to your research question and
discuss any uncertainty in your results. -- If appropriate, discuss why
the research question cannot be fully answered using the data and
techniques available.

\section{References}\label{references}

-- Cite data sources. -- Cite literature and other sources useful for
explaining the background informa- tion, research question, and
analytical tools.




\end{document}
